\begin{figure*}[tbhp]
    \centering

    \subfloat[Input]{
        \includegraphics[width=0.135\linewidth]{figs/SR/LR}
    }
    \subfloat[Ground Truth]{
        \includegraphics[width=0.135\linewidth]{figs/SR/HR}
    }
    \subfloat[Bicubic]{
        \includegraphics[width=0.135\linewidth]{figs/SR/Bicubic}
    }
    \subfloat[Baseline]{
        \includegraphics[width=0.135\linewidth]{figs/SR/Photography}
        \label{figure:comparisons/baseline}
    }
    \subfloat[Ours (MAE)]{
        \includegraphics[width=0.135\linewidth]{figs/SR/Ours@ESPCN@L1}
        \label{figure:comparisons/mae}
    }
    \subfloat[Ours (MSE)]{
        \includegraphics[width=0.135\linewidth]{figs/SR/Ours@ESPCN@L2}
        \label{figure:comparisons/mse}
    }
    \smallbreak
    \subfloat[Ours (DSSIM)]{
        \includegraphics[width=0.135\linewidth]{figs/SR/Ours@ESPCN@SSIM}
        \label{figure:comparisons/dssim}
    }
    \subfloat[Ours (MS-DSSIM)]{
        \includegraphics[width=0.135\linewidth]{figs/SR/Ours@ESPCN@MS-SSIM}
        \label{figure:comparisons/ms-dssim}
    }
    \subfloat[Ours (MixGE)]{
        \includegraphics[width=0.135\linewidth]{figs/SR/Ours@ESPCN@Sobel001}
        \label{figure:comparisons/mixge}
    }
    \subfloat[Ours (Sketch)]{
        \includegraphics[width=0.135\linewidth]{figs/SR/Ours@ESPCN@Sketch}
        \label{figure:comparisons/sketch}
    }
    \subfloat[Ours (Perceptual)]{
        \includegraphics[width=0.135\linewidth]{figs/SR/Ours@ESPCN@Perceptual}
        \label{figure:comparisons/perceptual}
    }

    \caption[]{
        Comparison across loss functions on the \textbf{Small} network.
        Regions outlined in magenta are magnified to better visualize the impact of applying different SISR methods. ``Ours (MixGE)'' uses $\lambda_G = 0.01$ and the Baseline is the pre-trained Y-channel ESPCN model.
        From top to bottom rows, images are cropped samples of the following sources:
        \begin{enumerate*}
            \item Manga109\cite{mtap_matsui_2017,multimedia_aizawa_2020}, \copyright~Ken~Akamatsu\label{figure:manga109}
            \item Wikipedia (\url{https://en.wikipedia.org/wiki/File:Wikipe-tan_face.svg})
            \item SYNLA dataset\cite{synla}
            \item Set14 dataset\cite{zeyde2010single}.
        \end{enumerate*}
    }

    \label{figure:comparisons}
\end{figure*}

